\paragraph{Proposed Approach.}
Markov semigroup methods provide a framework for proving influence
and isoperimetric inequalities \cite{CorderoErausquinLedoux12,ODonnellWimmer13,IvanisviliZhang26}.
Our work \cite{chen2026discrete} utilizes the similar idea together with gradient estimates.
We aim to apply the same framework to address the proposed problem.

A \emph{Markov semigroup} $(P_t)_{t\geq 0}$ describes the time evolution of a Markov process and is generated by an operator $L$ ($P_t = e^{tL}$): for any $f: \Omega\mapsto \R$ and $x\in\Omega$, $Lf(x) = \sum_{i=1}^n q_i(x)(f(x^i)-f(x))$.
Here $\{q_i\}_{i\in[n]}$ is called the \emph{transition rates} of the Markov semigroup, which is always assumed to be \emph{reversible} with respect to $\mu$: $\mu(x)q_i(x)=\mu(x^i)q_i(x^i)$.
Given the Markov semigroup $(P_t)_{t\geq 0}$, the corresponding strong gradient estimate holds with a constant $\rho$ if for any function $f:\Omega\mapsto\R$ and $t\geq 0$,
$$
    \norm{\nabla^\star P_t f}_2 \leq e^{-\rho t}P_t\tp{\norm{\nabla^\star f}_2},
$$ 
where the discrete gradient of $f$ is defined as $\nabla^\star f(x) = [q_i(x)(f(x^i)-f(x))]_{i\in[n]}$.
Through the potential function and derivative techniques, the gradient estimate can be equivalently written as a functional inequality, which is known as the reinforced Bakry--\'Emery condition \cite{BakryEmery85}.
This can be formulated as a local matrix eigenvalue problem \cite{CKLP22}.
In \cite{chen2026discrete}, the implication from the gradient estimates to isoperimetric inequalities was established. 
By analyzing commutators between the generator and discrete gradient, we derived sufficient local matrix criteria for gradient estimates and verified them under suitable weak-dependence assumptions on $\mu$.

In our previous work, we chose transition rates $q_i(x) = \frac 12 \sqrt{\mu(x^i)/\mu(x)}$ satisfying the square identity $q_i(x)q_j(x^i) = q_j(x)q_i(x^j)$ and used the commutativity of coordinate flips $(x^i)^j = (x^j)^i$ to establish the commutator identity $(L\nabla^\star - \nabla^\star L)f= H\cdot \nabla^\star f$, where $H: \Omega\mapsto\R^{n\times n}$.
However, for distributions supported on matroid independent sets, coordinate flips may leave the support and the commutativity of elements may not hold, making the square identity fail.
Establishing an appropriate commutator identity is therefore a central challenge.
One potential approach is to redefine gradients and  apply a model-based transition rates.
The directed-derivative identities for birth--death dynamics in
\cite{gobel2026simple} and the local Bochner framework in
\cite{chen2026spectral} suggest a possible route to derive a new commutator identity and establish the gradient estimate.